\section{Introduction}
\label{sec:introduction}

Remaining useful life (\RUL) estimation for lithium-ion cells is challenging
because degradation is multi-mechanistic and operating-condition dependent.
Purely data-driven models often struggle when long-horizon aging data is
limited, while purely mechanistic electrochemical models can be difficult to
parameterise for specific large-format cells without extensive
characterisation.

This work documents a pragmatic ``physics-guided surrogate'' approach:
\begin{enumerate}
  \item Use \pybamm \citep{Sulzer2021PyBaMM} as a physics prior to generate
  plausible trajectories within a validated operating envelope.
  \item Train a Neural ODE surrogate \citep{Chen2018NeuralODE} for \SOH dynamics
  as a function of cycle number and health features.
  \item Constrain the surrogate to be physically consistent at minimum by
  enforcing monotonic \SOH decrease.
\end{enumerate}

\paragraph{Scope limitations (must be explicit for defensibility).}
All tests in this repository are conducted at a single ambient temperature
(25\,$^\circ$C). Therefore, temperature is treated as constant and the study
does \emph{not} claim temperature generalisation. In addition, the dataset
currently contains a small number of cells, which constrains the extent to
which end-of-life predictions can be validated.

